\section{O}

% öffnen (to open)
\begin{verbentry}{
  style = {acc},
  toc = {öffnen (to open)},
  name = {öffnen},
  english = {to open},
  type = {regular, + Akkusativ},
  auxiliary = {haben}
}
 
    
     \vspace{5pt}

    \hrule
    
    \vspace{5pt}

  \begin{tabular}{rl}
     \multicolumn{2}{l}{\textbf{PRÄSENS} }\\
    ich & \textbf{öffne}\\
    du & \textbf{öffnest}\\
    er/sie/es & \textbf{öffnet}\\
    wir & \textbf{öffnen}\\
    ihr & \textbf{öffnet}\\
    sie/Sie & \textbf{öffnen}\\
  \end{tabular}
  \hfill
     \begin{tabular}{rl}
    \multicolumn{2}{l}{\textbf{PRÄTERITUM} }\\
    ich & \textbf{öffnete}\\
    du & \textbf{öffnetest}\\
    er/sie/es & \textbf{öffnete}\\
    wir & \textbf{öffneten}\\
    ihr & \textbf{öffnetet}\\
    sie/Sie & \textbf{öffneten}\\
  \end{tabular}
  \hfill
  \begin{tabular}{rl}
     \multicolumn{2}{l}{\textbf{IMPERATIV} }\\
   & \textbf{öffne (du)!}\\
   & \textbf{öffnen wir!}\\
   & \textbf{öffnet (ihr)!}\\
   & \textbf{öffnen Sie!}\\
   & \vspace{10pt}
  \end{tabular}

    \vspace{10pt}

 \begin{tabular}{rl}
     \multicolumn{2}{l}{\textbf{PERFEKT}}\\
   & haben + \textbf{geöffnet}\\
  \end{tabular}
\hfill
   \begin{tabular}{rl}
   
    \multicolumn{2}{l}{\textbf{FUTURE I} }\\
  &  werden + \textbf{öffnen}\\
 
  \end{tabular}
\hfill
   \begin{tabular}{rl}
      \multicolumn{2}{l}{\textbf{PLUSQUAMPERFEKT}}\\
  &  hatten + \textbf{geöffnet}\\
  \end{tabular}
  
  
\vspace{10pt}

\begin{tabular}{lll}
\multicolumn{3}{l}{\textbf{MIT PRÄPOSITIONEN}}\\
& \textbf{---} & \\
\end{tabular}
\hfill
\begin{tabular}{lll}
\multicolumn{3}{l}{\textbf{TRENNBARE VERBEN}}\\
& \textbf{---} & \\
\end{tabular}
\vspace{-5pt}
\end{verbentry}

% organisieren (to organize)
\begin{verbentry}{
  style = {acc},
  toc = {organisieren (to organize)},
  name = {organisieren},
  english = {to organize},
  type = {regular, + Akkusativ},
  auxiliary = {haben}
}


\vspace{5pt}
\hrule
\vspace{5pt}

\begin{tabular}{rl}
\multicolumn{2}{l}{\textbf{PRÄSENS}}\\
ich & \textbf{organisiere}\\
du & \textbf{organisierst}\\
er/sie/es & \textbf{organisiert}\\
wir & \textbf{organisieren}\\
ihr & \textbf{organisiert}\\
sie/Sie & \textbf{organisieren}\\
\end{tabular}
\hfill
\begin{tabular}{rl}
\multicolumn{2}{l}{\textbf{PRÄTERITUM}}\\
ich & \textbf{organisierte}\\
du & \textbf{organisiertest}\\
er/sie/es & \textbf{organisierte}\\
wir & \textbf{organisierten}\\
ihr & \textbf{organisiertet}\\
sie/Sie & \textbf{organisierten}\\
\end{tabular}
\hfill
\begin{tabular}{rl}
\multicolumn{2}{l}{\textbf{IMPERATIV}}\\
& \textbf{organisiere (du)!}\\
& \textbf{organisieren wir!}\\
& \textbf{organisiert (ihr)!}\\
& \textbf{organisieren Sie!}\\
\end{tabular}

\vspace{10pt}

\begin{tabular}{rl}
\multicolumn{2}{l}{\textbf{PERFEKT}}\\
& haben + \textbf{organisiert}\\
\end{tabular}
\hfill
\begin{tabular}{rl}
\multicolumn{2}{l}{\textbf{FUTURE I}}\\
& werden + \textbf{organisieren}\\
\end{tabular}
\hfill
\begin{tabular}{rl}
\multicolumn{2}{l}{\textbf{PLUSQUAMPERFEKT}}\\
& hatten + \textbf{organisiert}\\
\end{tabular}

\vspace{10pt}

\begin{tabular}{lll}
\multicolumn{3}{l}{\textbf{MIT PRÄPOSITIONEN}}\\
& \textbf{organisieren für} + Akk & (to organize for)\\
\end{tabular}
\hfill
\begin{tabular}{lll}
\multicolumn{3}{l}{\textbf{TRENNBARE VERBEN}}\\
& \textbf{---} & \\
\end{tabular}

\vspace{-5pt}
\end{verbentry}
